\documentclass[12pt]{book}
\usepackage{MLA_UPIITA}
\usepackage[scaled]{helvet}             % Tipografía 
\renewcommand\familydefault{\sfdefault} 
\usepackage[T1]{fontenc}
\titulo{Título}
\fechaentrega{\today}
\nombres{Nombre(s)}
\institucion{Institución}
\apellidos{Apellidos}
\profesor{Profesor}
\asignatura{Asignatura}
\fuentes{Bibliografia.bib}

\begin{document}
\pagenumbering{Roman}
\maketitlepage          %Portada
\newpage
\tableofcontents        %Índice general
\newpage
\listoffigures          %Índice de figuras
\newpage
\listoftables           %Índice de tablas

\newpage
\pagenumbering{arabic}
% Abstract
\begin{abstract}
	\lipsum[1-3]
\end{abstract}

% Para incluir encabezados basta con desmarcar la siguiente línea
%\makeheader 

% Insertar .tex al documento general
\newpage
\input{IMRaD/Introducción}

\newpage
\input{IMRaD/Métodos}

\newpage
\input{IMRaD/Resultados}

\newpage
\input{IMRaD/Conclusiones}

% Cita literal corta de 4 o menos líneas
% \citaliteralcorta{Texto de la cita}{Numero de página(s)}{Referencia al nombre de la fuente, por ejemplo: Libro1} 

% Cita literal larga de más de 4 líneas
% \citaliterallarga{Texto de la cita}{Numero de página(s)}{Referencia al nombre de la fuente, por ejemplo: Libro2} 

% Para hacer citas dentro del texto:
%\cite[Número de página(s)]{Referencia al nombre de la fuente, por ejemplo: Libro1}

% Niveles de indización
% \section{} ->       1
% \subsection{} ->    2
% \subsubsection{} -> 3
% \paragraph{} ->     4
% \subparagraph{} ->  5

\makeworkscited % Muestra las referencias en el documento
\end{document}
